\documentclass[]{article}

\usepackage[margin=2cm]{geometry}
\usepackage{setspace}
\usepackage{siunitx}
\usepackage[italian]{babel}
\usepackage{amsmath}
\usepackage{parskip}
\usepackage{cancel}
\usepackage{float}
\usepackage{amssymb}
\usepackage{mathtools}

\usepackage[colorlinks=true, linkcolor=blue, urlcolor=black]{hyperref}

\newcommand{\Lag}{\mathcal{L}}
\newcommand{\Ham}{\mathcal{H}}
\newcommand{\normord}[1]{:\! #1 \!:}

%opening
\author{Matti Siga}
\title{Appunti di fondamenti di teoria dei campi- WIP}
\date{A.A. 2025/2026}
\makeatletter \@addtoreset{section}{part} \makeatother

\begin{document}
\maketitle
\tableofcontents
\onehalfspacing

\section{Introduzione e richiami}
Si consideri la teoria dei campi libera di un campo scalare reale, la cui densità lagrangiana campo è: \footnote{Imponendo da ora dove non esplicitamente detto le unità naturali: $\hbar = c = 1$}
\begin{equation}
 \Lag = \frac{1}{2} \left(  \partial_\mu \phi^ 2 \right)  -\frac{1}{2} m^2 \phi^2 = +\frac{1}{2} \dot{\phi}^2 -\frac{1}{2} \left( \vec\nabla \phi \right)^2 - \frac{m^2}{2} \phi
\end{equation}
Si noti che in quanto l'azione è adimensionale allora $\left[\Lag \right] = M^4$ quindi il campo $\phi$ ha le dimensioni di una massa. Si possono ottenere le equazioni del moto classiche usando le equazioni di Eulero-Lagrange, ottenendo le equazioni di Klein - Gordon:
\begin{equation}
 \partial_\mu \frac{\partial \Lag}{\partial \partial_\mu \phi} - \frac{\partial \Lag}{\partial \phi} = 0 \Rightarrow \left( \square + m^2 \right) = 0
\end{equation}
Similarmente usando il campo coniugato $\pi \equiv \frac{\partial \Lag}{\partial \dot{\phi}} = \dot{\phi}$ si può ottenere la densità hamiltoniana del campo:
\begin{equation}
 \Ham = \pi \dot{\phi} - \Lag = \frac{1}{2} \pi^2 + \frac{1}{2} \left( \vec \nabla \phi \right)^2 + \frac{1}{2} m^2 \phi^2
\end{equation}
N.B. Si noti che $\Ham$ è definita positiva perchè il campo è reale e tutti i termini del campo sono al quadrato

Lavorando in analogia all'oscillatore armonico si quantizza il campo promovendo il campo scalare in un punto a operatore usando le regole di commutazione (sapendo però che si ha infiniti gradi di libertà). In Rappresentazione di Schrödinger (dove gli operatori sono indipendenti dal tempo) si ha che $\hat{\phi}(\vec{x})$ è l'espansione in onde piane del campo:
\begin{equation}
  \hat{\phi}(\vec{x}) = \int \frac{d^3 p}{(2 \pi)^3} \frac{1}{\sqrt{2E_p}} \left( \hat{a}_p e^{ i\, \vec p \cdot \vec x} + \hat{a}^\dagger_p e ^ {-i \, \vec p \cdot \vec x} \right)
\end{equation}

Dove $E_p$ è l'energia (o la frequenza) di una particella libera di massa $m$: $E_p = \sqrt{\hat{p}^2 + m^2}$. Similarmente il campo coniugato è:
\begin{equation}
  \hat{\pi}(\vec{x}) = \int \frac{d^3 p}{(2 \pi)^3} (-i) \sqrt{\frac{E_p}{2}} \left( \hat{a}_p e^{ i\, \vec p \cdot \vec x} - \hat{a}^\dagger_p e ^ {-i \, \vec p \cdot \vec x} \right)
\end{equation}
A questo punto si impongono le condizioni di quantizzazione:
\begin{equation}
\begin{cases}

   \left[ \hat{\phi}(\vec x) ; \: \hat{\pi} (\vec y) \right] = i \delta^{(3)} (x - y ) \\
   \\
   \left[ \hat{\phi}(\vec x) ; \: \hat{\phi} (\vec y) \right] = \Big[ \hat{\pi}(\vec x) ; \: \hat{\pi} (\vec y) \Big] = 0 \qquad \text{(campi in punti diversi non comunicano istantaneamente)}
\end{cases}
\end{equation}
Da queste si trova:
\begin{equation}
\begin{cases}

   \Big[ \hat{a}_{\vec p} ; \: \hat{a}_{\vec {p\prime}} ^ \dagger \Big] = (2\pi)^ 3 \delta^{(3)} (\vec p - \vec{p \prime} ) \\
   \\
   \Big[ \hat{a}_{\vec p} ; \: \hat{a}_{\vec {p\prime}} \Big] =  \left[ \hat{a}_{\vec p}^\dagger ; \: \hat{a}_{\vec {p\prime}} ^ \dagger \right] = 0
\end{cases}
\end{equation}
N.B. Si possono ricavare le stesse condizioni di quantizzazione imponendo anche queste ultime relazioni di commutazione ricavando le prime

L'hamiltoniana è quindi anch'essa un operatore ed è:
\begin{equation}
 \hat{H} = \int \frac{d^3 p}{(2\pi)^3} \, E_p \left( \hat{a}_{\vec{p}} \: \hat{a}_{\vec{p}}^\dagger +\frac{1}{2} \left[ \hat{a}_{\vec{p}}; \: \hat{a}_{\vec{p}}^\dagger \right]\right)
\end{equation}
Dove il commutatore non è nullo nella energia di punto zero, ma si ha un infinito in quanto vi è una Delta di Dirac, infatti considerando un numero infinito di oscillatori armonici nello stato di vuoto $| 0 \rangle $ l'energia ``esplode''. Si può togliere  il problema definendo $\hat{H}$ come:
\begin{equation}
 \hat{H} = \int d^3 x \, \normord{\Ham}
\end{equation}
Dove in $\normord{\Ham}$ viene eseguita l'operazione di normal ordering, dove vengono spostati gli operatori di distruzione a destra ($\normord{\hat{a}_p \, \hat{a}_p^\dagger} \; =  \hat{a}_p^\dagger  \, \hat{a}_p $), in modo tale che sullo stato di vuoto non si creano problemi in quanto $ \hat H |0\rangle = 0$ poiché $\hat a_p | 0 \rangle = 0$, mentre $|\vec p \rangle = \sqrt{2 E_p} \hat a ^ \dagger_p |0\rangle$ è la creazione di una particella di momento p (Questo è lo spazio di Fock, in cui l'impulso è definito in quanto gli operatori di creazione e distruzione sono costrutti rispetto ad esso).

Una delle possibili normalizzazioni dell'impulso è invece: $\langle \vec p | \vec q \rangle = 2 E_p (2 \pi)^3 \delta^{(3)} ( \vec p - \vec q)$, che ha il vantaggio di essere Lorentz invariante.

Ci si può chiedere cosa sia fisicamente $\hat \phi$. Dal momento che $\hat \phi |0\rangle$ fornisce una particella scalare nel punto x, $\hat \phi$ crea particelle in un punto definito. La dipendenza temporale del campo può essere espressa usando la Rappresentazione di Heisenberg (detto II quantizzazione):
\begin{equation}
 \hat \phi_H (t; \vec x) \equiv e^{i \hat H t} \hat \,  \phi(\vec x) \,  e^{-i \hat H t} = \hat \phi_H(x) \quad \text{dove x è un quadrivettore}\footnote{\text{Si può quindi distinguere un operatore scritto alla Heisenberg o alla Schrödinger a seconda che si usi $x$ o $\vec x$}}
\end{equation}
Si può inoltre osservare che con una espansione in onde piane si ha in tale rappresentazione:
\begin{equation}
 \hat \phi_H (x) = \int \frac{d^3 p}{(2\pi)^2} \frac{1}{\sqrt{2 E_p}}\left. \left( \hat a_{\vec{p}} e^{-ipx} + \hat a_{\vec{p}}^\dagger e^{+ipx} \right) \right|_{p^0 = E_p \enspace x^0 = t}
\end{equation}
Dove al primo termine sono associate frequenze positive e il secondo negative.\\
N.B. Non è l'unico modo possibile per fare la seconda quantizzazione, in quanto può essere fatta anche tramite integrali di cammino.

Si introduce inoltre il correlatore tra due campi come:
\begin{equation}
 D(x-y) \equiv \langle0| \hat \phi (x) \hat \phi (y) | 0 \rangle = \cdots = \int \frac{d^3 p}{(2 \pi)^3} \frac{1}{2 E_p} e^{- i p \cdot ( x -y)}
\end{equation}
Si ha però che esso non va a zero fuori dal conoluce dello spazio tempo, quindi non è necessariamente causale, ma lo è invece:
\begin{equation}
 \langle 0 | \left[ \hat \phi (x); \: \hat \phi (y) \right] | 0 \rangle = \cdots = D(x-y) - D(y -x)
\end{equation}
Infine si ricorda che il propagatore di Feymann è:
\begin{equation}
 D_F (x-y) = \langle 0 |T\hat \phi(y) \hat \phi(x) |0\rangle = \int \frac{d^3 p}{(2\pi)^3} \frac{i}{p^2 - m^2 +i\epsilon} e ^ {-i p \cdot (x-y)}
\end{equation}
Dove $T$ è il T-Prodotto così definito:
\begin{equation}
 T[\hat \phi(y) \hat \phi(x)] =
 \begin{cases}
  D(x-y) \quad \text{se } x^0 > y^0 \\
  D(y-x) \quad \text{se } y^0 > x^0
 \end{cases}
\end{equation}
E invece $\epsilon$ è un infinitesimo positivo tale per fare lungo l'asse reale $p_0$ in quanto vi sono dei poli in $p_0 = \pm(E_p - i\epsilon)$ (Vi sono altri cammini di integrazione possibili, ma questo è il più utile.)
\end{document}


